\section{Phân tích độ nhạy của phép đo sai số}

Phân vị 95 mô tả giá trị tại một vị trí trong phân bố,
còn max mô tả mẫu bất lợi nhất trong corpus. Chúng
phản ứng khác nhau khi một nhóm nhỏ giao dịch giữ
thêm số hạng. Từ n8 đến n9, số cặp lấy mẫu đều còn
term giảm từ 554 xuống 43; tới n10 bằng 0 trong nhóm
này. P95 gần với mức sai số do cắt đầu vào ngay từ
n8, nhưng max vẫn thay đổi theo những mẫu hiếm còn term.
Trong nhóm chủ đích, có cặp cần đủ 12 vòng. Vì vậy
không thể suy rằng mọi dữ liệu đều đủ ở n10.

\begin{engineeringtable}{Các mốc ở đuôi phân bố trên nhóm lấy đều, OPS1 (3 chữ số có nghĩa).}{tab:tail-sensitivity}
\begin{tabular}{rrrr}
\toprule
n & Cặp còn term & P95 sai số tổng (\%) & Max quan sát (\%) \\
\midrule
7 & 3.724 & 0,0320 & 0,351 \\
8 & 554 & 0,0299 & 0,156 \\
9 & 43 & 0,0297 & 0,0582 \\
10 & 0 & 0,0297 & 0,0460 \\
12 & 0 & 0,0297 & 0,0460 \\
\bottomrule
\end{tabular}
\end{engineeringtable}

Phân tầng regime giúp quan sát rõ các miền ít xuất hiện khi lấy đều theo mã. Với
$n=12$, nhóm phân tầng đạt $\mathrm{P95} = 0{,}0218\%$, giá trị cực đại $0{,}0438\%$
và trung bình $2{,}66$ vòng lặp (so với trung bình $5{,}03$ vòng ở nhóm lấy đều).
Sự khác biệt này phản ánh mật độ bit $1$ và độ dài fraction khác nhau giữa các dải
\textit{regime}; do đó số liệu của nhóm phân tầng được theo dõi tách biệt khỏi phân
bố ứng dụng thực tế.

\section{Từ sai số tích đến sai số MAC}

Sai số cắt fraction được bảo toàn giống nhau qua \texttt{MAC-NF} và \texttt{MAC-F}
cho cùng một tích đầu vào. Khi toán hạng cộng $C$ cùng dấu với $P$, độ lớn tổng tăng
nên cùng một sai lệch tích thường chiếm tỷ lệ tương đối nhỏ hơn; ngược lại, khi $C$
trái dấu, mẫu số thu hẹp và bước lượng tử hóa trung gian $R(\widehat P)$ của
\texttt{MAC-NF} làm xuất hiện sai khác phân kỳ so với tổng đầy đủ.

Trên $108.703$ ca gần $0$ ($r \ne 0$) thuộc nhóm ép triệt tiêu $C = R(-P)$, hai
chỉ số bất biến theo thang trong Bảng~\ref{tab:cancel-target-metrics} phản ánh bản
chất sai số chính xác hơn nhiều so với sai số tuyệt đối thô vốn biến thiên hàng
chục bậc theo hệ số mũ của \texttt{Posit(32,3)}. Đồng thời, vì $C$ bị lượng tử hóa
trước khi đưa vào RTL, một số cặp chạm ngưỡng kẹp hoặc bước nhảy thưa khiến $P+C$
không còn thỏa điều kiện gần $0$; việc phân lớp lại theo tổng chính xác $r = P + C$
trước khi thống kê giúp tách bạch hoàn toàn nhóm $r = 0$ khỏi các ca kẹp ngưỡng
ở lớp \texttt{regular}.

\section{Ý nghĩa của kết quả chuỗi}

Trong vận hành theo chuỗi, \texttt{MAC-NF} và \texttt{MAC-F} cập nhật trạng thái
\texttt{Posit(32,3)} sau mỗi phép hoàn tất nên sai khác làm tròn trung gian tích tụ
qua từng bước hồi tiếp. Trái lại, \texttt{MAC-Quire} tích lũy trực tiếp trên lưới
972 bit và chỉ làm tròn duy nhất tại lệnh \texttt{READ}. Ở nhóm chuỗi gần triệt tiêu
có nhiễu nhỏ tại bit thấp bị cắt của cặp toán hạng cuối, các tích đã cắt gọn triệt
tiêu lẫn nhau giúp \texttt{MAC-Quire} giữ sai số ở bậc $10^{-8}$, tránh lỗi tích tụ
trạng thái giữa chuỗi.

Tại độ dài $1.024$ MAC, độ trễ tuyệt đối theo xung nhịp thực (Bảng~\ref{tab:chain-cycles})
giúp phân định rõ ranh giới lựa chọn: khi đặt trong hệ thống xử lý thời gian thực
cần tần số cao và tài nguyên gọn, \texttt{MAC-F} nhanh hơn đáng kể so với cả hai
biến thể còn lại dù hai lõi \texttt{MAC-NF}/\texttt{MAC-F} có sai số chuỗi ngang nhau.

\begin{engineeringtable}{Sai số tuyệt đối cuối của nhóm gần triệt tiêu ($4$ chuỗi/lõi, 3 chữ số có nghĩa).}{tab:near-chain-detail}
\begin{tabular}{rrrrr}
\toprule
MAC/chuỗi & Trung vị NF/F & Max NF/F & Trung vị quire & Max quire \\
\midrule
16 & $6{,}94{\times}10^{-8}$ & $9{,}97{\times}10^{-8}$ & $9{,}47{\times}10^{-9}$ & $1{,}34{\times}10^{-8}$ \\
64 & $1{,}54{\times}10^{-7}$ & $7{,}07{\times}10^{-7}$ & $9{,}82{\times}10^{-9}$ & $1{,}02{\times}10^{-8}$ \\
256 & $2{,}14{\times}10^{-6}$ & $5{,}59{\times}10^{-6}$ & $8{,}30{\times}10^{-9}$ & $9{,}36{\times}10^{-9}$ \\
1.024 & $2{,}17{\times}10^{-5}$ & $1{,}28{\times}10^{-4}$ & $1{,}73{\times}10^{-8}$ & $2{,}61{\times}10^{-8}$ \\
\bottomrule
\end{tabular}
\end{engineeringtable}

Bảng~\ref{tab:near-chain-detail} dùng giá trị cuối so với
tổng tích đủ fraction, trên nhóm có các cặp tích đối dấu
và một perturbation ở bit thấp của cặp cuối. Với bốn mẫu
mỗi độ dài, các hàng mô tả giá trị đã quan sát. Chúng không
đủ để suy một luật tăng sai số theo chiều dài cho mọi dãy.
Giá trị trung vị của NF và F bằng nhau trong nhóm này,
dù số chu kỳ khác nhau. Vì vậy lựa chọn theo tài nguyên,
chu kỳ và độ chính xác phải giữ ba đại lượng tách biệt.

\section{Phân tích đường găng và tối ưu hóa vật lý sau định tuyến}

\subsection{Phân biệt tài nguyên giữa tổng hợp và định tuyến}

Ở pha tổng hợp cơ sở (Bảng~\ref{tab:nf-f-synthesis}), \texttt{MAC-F} dùng ít
hơn 21,9\% LUT (347 LUT) và 14,7\% FF (135 FF) so với \texttt{MAC-NF}. Chênh
lệch FF xuất phát trực tiếp từ ngân hàng thanh ghi của khối đóng gói tích trung
gian mà \texttt{MAC-F} loại bỏ, còn tỷ lệ LUT phụ thuộc thêm vào chiến lược gộp
logic điều khiển và ánh xạ chuỗi cộng của Vivado qua từng pha triển khai.

Sau định tuyến OOC 100\,MHz, cấu hình trước tách tải fanout của \texttt{MAC-F}
dùng 1.211 LUT / 781 FF (so với 1.558 LUT / 916 FF của \texttt{MAC-NF}). Khi
tách bốn nhánh cho tín hiệu hoán đổi toán hạng tại tầng chuẩn bị S1, thiết kế
hoàn thiện của \textbf{MAC-F} sử dụng 1.221 LUT (chỉ tăng $+10$ LUT) và giữ nguyên
781 FF, trong khi \texttt{MAC-NF} dùng 1.591 LUT và 916 FF; như vậy \textbf{MAC-F}
tiếp tục tiết kiệm 370 LUT (\textbf{23,3\%}) và 135 FF (\textbf{14,7\%}) so với \texttt{MAC-NF}.

\subsection{Phân tích sụt giảm timing hậu route và hiệu quả tách fanout tầng S1}

\begin{engineeringtable}{Các đường găng quyết định slack sau định tuyến OOC 100\,MHz: tách riêng đường cổng \texttt{in\_ready} và đường \textit{reg-to-reg} nội bộ trước và sau tối ưu tầng S1.}{tab:route-critical}
\begin{tabularx}{\linewidth}{@{}>{\raggedright\arraybackslash}p{34mm}>{\raggedright\arraybackslash}p{15mm}Y Y@{}}
\toprule
Lõi (trạng thái) & Slack (ns) & Điểm \textit{launch} & Điểm \textit{capture} \\
\midrule
NF setup (I/O, trước tách) & $+0{,}665$ & Packer cuối, \texttt{meta\_sync2\_reg} & Cổng \texttt{in\_ready}. \\
NF setup (R2R, trước tách) & $+2{,}943$ & \texttt{u\_product\_pack/\allowbreak d\_q\_reg[15]} & \texttt{product\_frac\_\allowbreak q\_reg[5]} (Parser tích). \\
NF setup (I/O, tách S1) & $+0{,}669$ & \texttt{u\_pack/p1\_valid\_\allowbreak q\_reg} & Cổng \texttt{in\_ready}. \\
NF setup (R2R, tách S1) & $+2{,}639$ & \texttt{u\_core/acc\_q\_reg[20]} & \texttt{u\_product\_pack/\allowbreak p1\_guard\_bit\_q\_reg} (\texttt{dp=7,256\,ns}). \\
F setup (I/O, trước tách) & $+2{,}142$ & Wrapper, \texttt{meta\_sync2\_reg} & Cổng \texttt{in\_ready}. \\
F setup (R2R, trước tách) & $+2{,}768$ & \texttt{core\_sf\_\allowbreak inc\_q\_reg[4]} & \texttt{u\_add/\allowbreak s1\_sf\_q\_reg[1]} (Adder S1). \\
\textbf{F setup (I/O, hoàn thiện)} & \textbf{$+2{,}247$} & \texttt{u\_b/p2\_valid\_q\_reg} & Cổng \texttt{in\_ready}. \\
\textbf{F setup (R2R, hoàn thiện)} & \textbf{$+3{,}524$} & \texttt{c\_frac\_q\_reg[9]} & \texttt{u\_add/\allowbreak s1\_x\_frac\_q\_reg[1]} (\texttt{dp=6,321\,ns}). \\
Quire setup (R2R) & $-0{,}187$ & \texttt{u\_bank/\allowbreak offset\_q\_reg[4]} & \texttt{g\_words[16].\allowbreak q\_words\_reg[16][8]/CE}. \\
NF hold (R2R, tách S1) & $+0{,}041$ & \texttt{u\_product\_pack/\allowbreak p1\_special\_sel\_q\_reg[2]} & \texttt{u\_product\_pack/\allowbreak d\_q\_reg[27]}. \\
F hold (R2R, hoàn thiện) & $+0{,}043$ & Adder, \texttt{flags\_out\_reg[0]} & Packer P1, \texttt{p1\_flags\_q\_reg[0]}. \\
Quire hold (R2R) & $+0{,}040$ & \texttt{pending\_words\_\allowbreak reg[11][32]} & \texttt{q\_words\_\allowbreak reg[11][32]}. \\
\bottomrule
\end{tabularx}
\end{engineeringtable}

Độ trễ F2F trong Bảng~\ref{tab:f2f-stages} (Phụ lục~\ref{app:fpga}) phản ánh
thời gian lan truyền tổ hợp lớn nhất tới từng nhóm thanh ghi đích sau tổng hợp,
dùng để nhận diện sớm tầng cần theo dõi. WNS toàn thiết kế sau định tuyến bao
gồm thêm độ lệch clock và quỹ định thời đường cổng ngoài (chi tiết từng thanh ghi
được ghi nhận tại Bảng~\ref{tab:route-critical}). Dữ liệu hậu \textit{route}
làm rõ hiện tượng đảo chiều định thời ở cấu hình trước tối ưu và hiệu quả của kỹ
thuật tách tải:
\begin{itemize}
  \item Ở pha tổng hợp, \texttt{MAC-F} vượt trội hoàn toàn với WNS \textit{reg-to-reg}
        đạt $+4{,}319$\,ns (5 tầng logic) so với $+3{,}435$\,ns của \texttt{MAC-NF}
        (6 tầng logic).
  \item Tuy nhiên sau định tuyến, đường găng nội bộ \textit{reg-to-reg} của
        \texttt{MAC-F} ban đầu sụt từ $+4{,}319$\,ns xuống còn $+2{,}768$\,ns
        ($138{,}27$\,MHz, trễ đường dữ liệu $7{,}018$\,ns với \textbf{73,67\% là
        trễ định tuyến}). Nguyên nhân vật lý nằm ở nhánh chuẩn hóa từ thanh ghi
        số mũ tích đi vào khối tiền xử lý cộng, nơi tín hiệu chọn toán hạng có
        độ phân tải tới $62$ nhánh song song trên toàn độ rộng 26 bit fraction và
        metadata tầng S1 (cùng nút tải $130$ nhánh ở bộ tích lũy), gây nghẽn
        định tuyến cục bộ.
\end{itemize}

Giải pháp nhân bản bốn nhánh cân bằng tải tại khối chuẩn bị S1 (giữ nguyên toàn
bộ datapath, số FF và độ trễ chu kỳ) đã dịch chuyển đường găng nội bộ sang nhánh
căn chỉnh thanh ghi cộng với tỷ lệ trễ cân đối hơn (38,78\% logic / 61,22\% định
tuyến), nâng WNS \textit{reg-to-reg} của \textbf{MAC-F} lên \textbf{$+3{,}524$\,ns}
($F_{\max,\mathrm{R2R}} = 154{,}42$\,MHz, tăng $+16{,}15$\,MHz) và WNS toàn thiết
kế đạt $+2{,}247$\,ns (tương ứng $F_{\max}$ toàn thiết kế $= 129{,}0$\,MHz trên
đường tổ hợp ra \texttt{in\_ready}). Trên cùng khối S1 dùng chung, \texttt{MAC-NF}
đạt WNS toàn thiết kế $+0{,}669$\,ns ($107{,}2$\,MHz) và WNS nội bộ $+2{,}639$\,ns
($135{,}85$\,MHz, giảm nhẹ $0{,}304$\,ns so với trước tách S1 do dao động đặt--định
tuyến ngẫu nhiên của một lần chạy P\&R đơn lẻ khi đường găng chuyển về nhánh tích
đến \texttt{u\_product\_pack} chứ không phải nghẽn tại S1; ngược lại, mức cải thiện
$+0{,}756$\,ns của \texttt{MAC-F} gắn trực tiếp với việc giải tỏa đường găng S1).
Như vậy, \textbf{MAC-F} vượt trội đồng thời ở tần số nội bộ ($+18{,}57$\,MHz),
tần số toàn thiết kế ($+21{,}8$\,MHz), độ trễ xử lý ($-3$ chu kỳ/MAC) và tài nguyên LUT/FF.

\section{Phân tích kết quả triển khai MAC-Quire}

Quá trình tối ưu \texttt{MAC-Quire} gồm hai pha rõ rệt. Thứ nhất, chuyển đổi kiến
trúc backend từ dịch toàn vectơ 972 bit sang cập nhật cục bộ theo từng từ 36 bit
giúp cắt giảm 71,6\% LUT (10.156 xuống 2.888), 40,9\% FF (4.458 xuống 2.636) và
gần một nửa độ trễ F2F sau tổng hợp. Thứ hai, khi định tuyến đầy đủ OOC 100\,MHz
(4.594 net sạch, 0 lỗi DRC nghiêm trọng), mẫu \texttt{MAC-Quire} đạt $+0{,}170$\,ns
sau tổng hợp nhưng chuyển sang vi phạm setup nhẹ $-0{,}187$\,ns sau định tuyến
($\mathrm{TNS}=-19{,}336$\,ns, $F_{\max,\mathrm{R2R}} = 98{,}16$\,MHz; hold đạt
$+0{,}040$\,ns). Đường găng dài $9{,}851$\,ns (17 tầng logic, gồm 9 tầng chuỗi cộng)
rơi vào mạch giải mã địa chỉ từ tới chân cho phép của ngân hàng thanh ghi quire
(Bảng~\ref{tab:route-critical}). Như vậy, khi cần triển khai IP chạy trên $100$\,MHz
với tài nguyên gọn (ít hơn 2,35 lần LUT và 3,38 lần FF), \textbf{MAC-F} là lựa chọn
tối ưu hơn, trong khi \texttt{MAC-Quire} phù hợp cho ứng dụng triệt tiêu dải rộng
chấp nhận xung nhịp thấp hơn hoặc chèn thêm tầng thanh ghi ở đường chọn từ.
